\section{Planteamiento del problema}
\label{sec:planteamiento}

El plan separa \textbf{tres niveles} de problema dentro de la investigación basada en diseño. No
son tres redacciones del mismo enunciado: cada nivel encadena con el siguiente hasta llegar al
problema que el artefacto debe resolver.

\subsection{Situación real (problemática)}

Es el §1.2: síntomas del sector bancario ---fragmentación, costos de integración, casos como TSB o
\textit{open banking} con APIs heterogéneas--- descritos con hechos y cifras. Todavía no es el
objeto de estudio ni la fila de problema de la matriz.

\subsection{Situación subyacente (problema abstracto)}

Al quitar el contexto del banco concreto permanecen dos \textbf{estructuras jerárquicas}
comparables: el \textbf{contrato OpenAPI} (\textit{paths}, \textit{operations}, \textit{schemas}) y
el \textbf{\ServiceDomain{} BIAN} (\Behavior{}, \BusinessObject{}, relaciones de control).

El diagnóstico es de \textbf{correspondencia}: mismos conceptos con nombres distintos, operaciones
HTTP que no reflejan \textit{behaviors}, esquemas que no cubren \textit{business objects},
versiones sin trazabilidad (§\ref{sec:problematica}; Tabla~\ref{tab:tipologias}).

En este nivel se fija el \textbf{objeto de estudio} (contrato OpenAPI bancario completo) y la
\textbf{unidad de análisis} (pares endpoint/schema $\leftrightarrow$ \Behavior{}/\BusinessObject{};
§\ref{sec:objeto-estudio}). BIAN funciona como \textbf{modelo de referencia} documental frente a la
capa REST.

\subsection{Problema tecnológico (situación a modificar)}

El \textbf{problema tecnológico} ---el de la matriz de consistencia y del objetivo general del
Capítulo~\ref{cap:objetivos}--- no coincide con la problemática del sector ni con la brecha
abstracta. Existen prácticas aisladas de revisión manual y herramientas genéricas de
\textit{schema matching}, pero \textbf{no hay un método reproducible y adaptado} que, dado un
contrato OpenAPI bancario y un \ServiceDomain{} BIAN emparejado, calcule de forma sistemática en
qué medida se alinean y bajo qué tipologías se desalinea. Hoy esa evaluación suele ser manual,
parcial o simplemente no ocurre antes del despliegue.

Este plan responde en el nivel de la \textbf{necesidad del artefacto}: diseño del
procedimiento y del modelo S/E/C $\rightarrow$ \AlignmentScore{}. Operar el artefacto en producción
bancaria queda fuera de alcance (§1.6).

\subsection{Problema de investigación}

Con los niveles anteriores en mente, el \textbf{problema de investigación} queda como pregunta
medible, distinta de los objetivos:

\begin{quote}
  \textbf{¿En qué medida y bajo qué tipologías se produce desalineación entre contratos OpenAPI y
  los modelos de negocio BIAN en arquitecturas API bancarias, y cómo cuantificarla de forma
  reproducible mediante un modelo de alineación?}
\end{quote}

Esta pregunta ancla el \textbf{problema tecnológico}; la brecha del estado del arte (Cap.~3 §3.6)
explica por qué la literatura aún no lo resuelve, sin confundirse con la formulación del problema
tecnológico en Anexo~C.

Para operacionalizar el problema se plantean las siguientes \textbf{preguntas específicas}:

\begin{table}[H]
  \centering
  \caption{Preguntas específicas de investigación}
  \label{tab:preguntas}
  \small
  \begin{tabular}{@{}p{0.06\textwidth}p{0.88\textwidth}@{}}
    \toprule
    \textbf{ID} & \textbf{Pregunta} \\
    \midrule
    \textbf{P1} &
    ¿Qué correspondencias estructurales pueden establecerse entre elementos de un contrato OpenAPI
    (\textit{paths}, \textit{operations}, \textit{schemas}) y los artefactos de un \ServiceDomain{}
    BIAN? \\
    \textbf{P2} &
    ¿Qué grado de similitud semántica existe entre los términos y conceptos de negocio en BIAN y los
    nombres de recursos, propiedades y descripciones en OpenAPI? \\
    \textbf{P3} &
    ¿Qué tipologías de desalineación predominan al comparar ambos artefactos en un mismo dominio de
    servicio? \\
    \textbf{P4} &
    ¿Es posible integrar métricas estructurales y semánticas en un indicador de coherencia global
    por cada instancia evaluada (par contrato OpenAPI bancario --- \ServiceDomain{} BIAN)? \\
    \bottomrule
  \end{tabular}
\end{table}

La \Cref{tab:tipologias} resume la tipología de desalineaciones identificadas:

\begin{table}[H]
  \centering
  \caption{Tipología de desalineaciones OpenAPI--BIAN}
  \label{tab:tipologias}
  \begin{tabular}{@{}p{0.18\textwidth}p{0.38\textwidth}p{0.32\textwidth}@{}}
    \toprule
    \textbf{Tipo} & \textbf{Descripción} & \textbf{Ejemplo ilustrativo} \\
    \midrule
    \textbf{Estructural} &
    Divergencia entre la organización de endpoints/esquemas y la estructura del SD BIAN &
    Operación \texttt{POST /transfer} sin correspondencia con el \Behavior{} \texttt{Initiate} definido en BIAN \\
    \textbf{Semántica} &
    Equivalencia estructural parcial pero significado de negocio distinto &
    Propiedad \texttt{amount} sin vínculo explícito con \texttt{PaymentAmount} del modelo BIAN \\
    \textbf{De modelado} &
    Representación incorrecta o incompleta de entidades de negocio &
    Esquema que omite atributos obligatorios del \BusinessObject{} BIAN \\
    \textbf{REST / contrato} &
    Inconsistencias en el estilo REST que dificultan la interpretación alineada al dominio &
    Uso de verbos en URLs o recursos que no mapean a \textit{Service Operations} BIAN \\
    \bottomrule
  \end{tabular}
\end{table}


Las preguntas P1--P4 orientan la investigación; las transformaciones concretas (normalización,
\textit{matching}, cálculo de \textit{scores}) pasan a los \textbf{objetivos específicos} del
Capítulo~\ref{cap:objetivos} y a las cinco fases declaradas en Anexo~C.
